% =====================================================================
\lezione{30/10/2026}{Laboratorio 3: risk register ISO/IEC 27001}
  {Gianluca Gilardi}
  {compilazione in gruppo di un risk register ISO/IEC~27001:2022 (controlli tecnologici dell'Allegato~A,
   clausola~8) per uno studio professionale}
\label{lez:2026-10-30}
% =====================================================================

\begin{nota}[Fonte]
Trascrizione automatica della registrazione: quasi tutta silenzio (gruppi al lavoro). Contenuti ripresi
dal briefing, dal playbook e dal file Excel del laboratorio (\texttt{lab/lab\_03/materiale}). Il caso,
il metodo e uno svolgimento ragionato sono nel Capitolo~\ref{lab:03}.
\end{nota}

\section{Sintesi}
Laboratorio in gruppo di circa 90 minuti. Ogni gruppo, nel ruolo di responsabile della sicurezza delle
informazioni di uno studio professionale, completa la valutazione e il trattamento di 18 scenari di
rischio già individuati: probabilità e impatto inerenti per riservatezza, integrità e disponibilità,
controlli dell'Allegato~A (solo clausola~8), misure, giustificazione, efficacia, rischio residuo e
decisione di trattamento.

\section{Organizzazione}
\begin{itemize}
  \item Il file dell'esercitazione è stato caricato nella cartella dei materiali della piattaforma poco
    prima dell'inizio; chi aveva già provato la copia di prova e letto il playbook partiva avvantaggiato.
  \item Prima cosa da compilare: nome del gruppo e tutti i partecipanti nel primo foglio, altrimenti il
    lavoro non può essere attribuito né valutato. Nome del file e caricamento secondo le istruzioni del
    file (come nel laboratorio precedente).
  \item \textbf{Consegna entro le 13:15}: i file caricati dopo vengono scartati.
  \item Esiti promessi entro domenica; i gruppi selezionati presentano lunedì 02/11.
\end{itemize}

\section{Contenuti: le indicazioni date in apertura}
Gianluca ripercorre il foglio di lavoro:
\begin{itemize}
  \item i rischi, le minacce e gli asset sono già compilati: si inseriscono probabilità e impatto
    inerenti, mentre il punteggio è calcolato in automatico dal file;
  \item si indicano i controlli applicabili, le misure concrete e il \textbf{perché} di quelle misure,
    e si valuta quanto sono efficaci rispetto a quel rischio specifico;
  \item si stimano probabilità e impatti residui (C, I, A) e si sceglie la decisione di trattamento
    (ridotto, accettato, trasferito, evitato);
  \item la propensione al rischio non ammette nulla sopra il livello medio: un rischio alto non si può
    semplicemente accettare;
  \item il foglio riepilogativo verifica solo la completezza, non la correttezza.
\end{itemize}

\begin{esame}
``Non esiste una risposta corretta'': tutto può essere discusso, purché il ragionamento sia ragionevole e
giustificato. Il playbook suggerisce una soluzione di riferimento, ma non è vincolante. Lo stesso criterio
vale per l'esame: conta la motivazione delle scelte di rischio inerente, controlli e rischio residuo.
\end{esame}

\section{Domande aperte e punti da verificare}
\begin{verifica}
\begin{itemize}
  \item Al termine Gianluca accenna a un intervento del professore, non registrato in modo comprensibile.
  \item Criteri di valutazione e soglia per le presentazioni: vedere gli avvisi in piattaforma.
\end{itemize}
\end{verifica}
